\newcommand{\jqasabstractbody}{%
In rugby union and association football, team key performance indicators (KPIs) can be represented in absolute terms or relative to the opponent. Relativisation sometimes improves match-outcome prediction and sometimes harms it. There has been no general account of when each occurs.\par
The answer depends on how much the two teams differ in variability and how strongly their KPI values rise and fall together. We combine these properties into the Paired Efficiency Factor (PEF), which generalises Fisher's paired-efficiency result to the unequal-variance conditions typical of competitive sport. Combined KPIs can interact in complex ways, so we analyse each indicator on its own.\par
Across $\PEFtotalStudies$ team KPIs from professional rugby union and association football, the PEF places every metric in one of four regimes. Anti-correlation is common, especially for high-volume competitive counts, and absolute measures are then usually preferable. Relativisation can still improve prediction when a noisier difference carries more outcome information. An idealised simulation and one representative KPI from each regime confirm both signs under team-blocked cross-validation.\par
The PEF turns an ad hoc feature-engineering choice into a data-informed diagnostic for when to relativise performance metrics and when not to.}
